% ========================================================
% BÀI TẬP TOÁN 5 - KẾT NỐI TRI THỨC VỚI CUỘC SỐNG
% Bài 39: Tìm hai số khi biết hiệu và tỉ số
% Định dạng: BTTeX Part III - Trắc nghiệm điền số ngắn
% ========================================================

% Câu 1
\begin{ex}
  Hai số có hiệu là 20 và tỉ số là $\frac{1}{2}$. Tìm số bé.
  \shortans{20}
  \loigiai{
    Hiệu số phần bằng nhau là 2 - 1 = 1 phần. Giá trị một phần là 20 : 1 = 20. Số bé là 20.
  }
\end{ex}

% Câu 2
\begin{ex}
  Hai số có hiệu là 20 và tỉ số là $\frac{1}{2}$. Tìm số lớn.
  \shortans{40}
  \loigiai{
    Hiệu số phần bằng nhau là 2 - 1 = 1 phần. Giá trị một phần là 20 : 1 = 20. Số lớn là 40.
  }
\end{ex}

% Câu 3
\begin{ex}
  Hai số có hiệu là 30 và tỉ số là $\frac{2}{5}$. Tìm số bé.
  \shortans{20}
  \loigiai{
    Hiệu số phần bằng nhau là 5 - 2 = 3 phần. Giá trị một phần là 30 : 3 = 10. Số bé là 20.
  }
\end{ex}

% Câu 4
\begin{ex}
  Hai số có hiệu là 30 và tỉ số là $\frac{2}{5}$. Tìm số lớn.
  \shortans{50}
  \loigiai{
    Hiệu số phần bằng nhau là 5 - 2 = 3 phần. Giá trị một phần là 30 : 3 = 10. Số lớn là 50.
  }
\end{ex}

% Câu 5
\begin{ex}
  Hai số có hiệu là 45 và tỉ số là $\frac{1}{4}$. Tìm số bé.
  \shortans{15}
  \loigiai{
    Hiệu số phần bằng nhau là 4 - 1 = 3 phần. Giá trị một phần là 45 : 3 = 15. Số bé là 15.
  }
\end{ex}

% Câu 6
\begin{ex}
  Hai số có hiệu là 45 và tỉ số là $\frac{1}{4}$. Tìm số lớn.
  \shortans{60}
  \loigiai{
    Hiệu số phần bằng nhau là 4 - 1 = 3 phần. Giá trị một phần là 45 : 3 = 15. Số lớn là 60.
  }
\end{ex}

% Câu 7
\begin{ex}
  Hai số có hiệu là 24 và tỉ số là $\frac{3}{5}$. Tìm số bé.
  \shortans{36}
  \loigiai{
    Hiệu số phần bằng nhau là 5 - 3 = 2 phần. Giá trị một phần là 24 : 2 = 12. Số bé là 36.
  }
\end{ex}

% Câu 8
\begin{ex}
  Hai số có hiệu là 24 và tỉ số là $\frac{3}{5}$. Tìm số lớn.
  \shortans{60}
  \loigiai{
    Hiệu số phần bằng nhau là 5 - 3 = 2 phần. Giá trị một phần là 24 : 2 = 12. Số lớn là 60.
  }
\end{ex}

% Câu 9
\begin{ex}
  Hai số có hiệu là 40 và tỉ số là $\frac{1}{3}$. Tìm số bé.
  \shortans{20}
  \loigiai{
    Hiệu số phần bằng nhau là 3 - 1 = 2 phần. Giá trị một phần là 40 : 2 = 20. Số bé là 20.
  }
\end{ex}

% Câu 10
\begin{ex}
  Hai số có hiệu là 40 và tỉ số là $\frac{1}{3}$. Tìm số lớn.
  \shortans{60}
  \loigiai{
    Hiệu số phần bằng nhau là 3 - 1 = 2 phần. Giá trị một phần là 40 : 2 = 20. Số lớn là 60.
  }
\end{ex}

% Câu 11
\begin{ex}
  Hai số có hiệu là 18 và tỉ số là $\frac{4}{7}$. Tìm số bé.
  \shortans{24}
  \loigiai{
    Hiệu số phần bằng nhau là 7 - 4 = 3 phần. Giá trị một phần là 18 : 3 = 6. Số bé là 24.
  }
\end{ex}

% Câu 12
\begin{ex}
  Hai số có hiệu là 18 và tỉ số là $\frac{4}{7}$. Tìm số lớn.
  \shortans{42}
  \loigiai{
    Hiệu số phần bằng nhau là 7 - 4 = 3 phần. Giá trị một phần là 18 : 3 = 6. Số lớn là 42.
  }
\end{ex}

% Câu 13
\begin{ex}
  Hai số có hiệu là 35 và tỉ số là $\frac{2}{7}$. Tìm số bé.
  \shortans{14}
  \loigiai{
    Hiệu số phần bằng nhau là 7 - 2 = 5 phần. Giá trị một phần là 35 : 5 = 7. Số bé là 14.
  }
\end{ex}

% Câu 14
\begin{ex}
  Hai số có hiệu là 35 và tỉ số là $\frac{2}{7}$. Tìm số lớn.
  \shortans{49}
  \loigiai{
    Hiệu số phần bằng nhau là 7 - 2 = 5 phần. Giá trị một phần là 35 : 5 = 7. Số lớn là 49.
  }
\end{ex}

% Câu 15
\begin{ex}
  Hai số có hiệu là 50 và tỉ số là $\frac{1}{6}$. Tìm số bé.
  \shortans{10}
  \loigiai{
    Hiệu số phần bằng nhau là 6 - 1 = 5 phần. Giá trị một phần là 50 : 5 = 10. Số bé là 10.
  }
\end{ex}

% Câu 16
\begin{ex}
  Hai số có hiệu là 50 và tỉ số là $\frac{1}{6}$. Tìm số lớn.
  \shortans{60}
  \loigiai{
    Hiệu số phần bằng nhau là 6 - 1 = 5 phần. Giá trị một phần là 50 : 5 = 10. Số lớn là 60.
  }
\end{ex}

% Câu 17
\begin{ex}
  Hai số có hiệu là 28 và tỉ số là $\frac{3}{7}$. Tìm số bé.
  \shortans{21}
  \loigiai{
    Hiệu số phần bằng nhau là 7 - 3 = 4 phần. Giá trị một phần là 28 : 4 = 7. Số bé là 21.
  }
\end{ex}

% Câu 18
\begin{ex}
  Hai số có hiệu là 28 và tỉ số là $\frac{3}{7}$. Tìm số lớn.
  \shortans{49}
  \loigiai{
    Hiệu số phần bằng nhau là 7 - 3 = 4 phần. Giá trị một phần là 28 : 4 = 7. Số lớn là 49.
  }
\end{ex}

% Câu 19
\begin{ex}
  Hai số có hiệu là 60 và tỉ số là $\frac{2}{5}$. Tìm số bé.
  \shortans{40}
  \loigiai{
    Hiệu số phần bằng nhau là 5 - 2 = 3 phần. Giá trị một phần là 60 : 3 = 20. Số bé là 40.
  }
\end{ex}

% Câu 20
\begin{ex}
  Hai số có hiệu là 60 và tỉ số là $\frac{2}{5}$. Tìm số lớn.
  \shortans{100}
  \loigiai{
    Hiệu số phần bằng nhau là 5 - 2 = 3 phần. Giá trị một phần là 60 : 3 = 20. Số lớn là 100.
  }
\end{ex}

% Câu 21
\begin{ex}
  Hai số có hiệu là 15 và tỉ số là $\frac{1}{2}$. Tìm số bé.
  \shortans{15}
  \loigiai{
    Hiệu số phần bằng nhau là 2 - 1 = 1 phần. Giá trị một phần là 15 : 1 = 15. Số bé là 15.
  }
\end{ex}

% Câu 22
\begin{ex}
  Hai số có hiệu là 15 và tỉ số là $\frac{1}{2}$. Tìm số lớn.
  \shortans{30}
  \loigiai{
    Hiệu số phần bằng nhau là 2 - 1 = 1 phần. Giá trị một phần là 15 : 1 = 15. Số lớn là 30.
  }
\end{ex}

% Câu 23
\begin{ex}
  Hai số có hiệu là 36 và tỉ số là $\frac{1}{4}$. Tìm số bé.
  \shortans{12}
  \loigiai{
    Hiệu số phần bằng nhau là 4 - 1 = 3 phần. Giá trị một phần là 36 : 3 = 12. Số bé là 12.
  }
\end{ex}

% Câu 24
\begin{ex}
  Hai số có hiệu là 36 và tỉ số là $\frac{1}{4}$. Tìm số lớn.
  \shortans{48}
  \loigiai{
    Hiệu số phần bằng nhau là 4 - 1 = 3 phần. Giá trị một phần là 36 : 3 = 12. Số lớn là 48.
  }
\end{ex}

% Câu 25
\begin{ex}
  Hai số có hiệu là 42 và tỉ số là $\frac{2}{9}$. Tìm số bé.
  \shortans{12}
  \loigiai{
    Hiệu số phần bằng nhau là 9 - 2 = 7 phần. Giá trị một phần là 42 : 7 = 6. Số bé là 12.
  }
\end{ex}
